\documentclass[
  12pt,
  openright,
  a4paper,
  twoside,
  english,
  brazil,
]{article}

\usepackage{mathtools}
\usepackage{amsfonts}
\usepackage{amssymb}
\usepackage{amsmath}
\usepackage[utf8]{inputenc}
\usepackage[portuguese]{babel}
\usepackage[T1]{fontenc}

\usepackage[top=2cm, bottom=2cm, left=2cm, right=2cm]{geometry}
\usepackage{enumerate}

\newcommand{\cqd}{\hfill $\square$}

% -----------------------------------------------------------------------------

\title{Teoria dos Números: Lista 01}
\author{GIPMAT 2026 - Zoeh}
\date{}

\begin{document}

\maketitle

\begin{enumerate}[\bfseries  1.]

  \item Prove por indução.

    Seja $a, b\in A$ e $m, n\in\mathbb N^*$. Então:

    \begin{enumerate}[\bfseries (a)]
      \item $(a^m)^n=a^{m\cdot n}$

        \textbf{Solução:}

        Seja $m$ um natural fixo e $S=\{n\in\mathbb N:(a^m)^n=a^{m\cdot n}\}$

        \begin{enumerate}[1.]
          \item $(a^m)^1= a^m=a^{m\cdot 1}$. Portanto $1\in S$
          \item Considere $k\in S$, de modo que $(a^m)^k=a^{m\cdot k}$

            Podemos multiplicar ambos os lados por $a^m$:

            $$
            (a^m)^k\cdot (a^m)=a^{m\cdot k}\cdot a^m
            $$
            $$
            (a^m)^{k+1}=a^{m\cdot k+m}=a^{m(k+1)}\implies
            (a^m)^{s(k)}=a^{m\cdot s(k)}
            $$

            Logo, $s(k)\in S$. E pelo Princípio da Indução Finita a
            propriedade é verdadeira para todos os números naturais.

        \end{enumerate}

        Agora considere $n$ um natural fixo e $S=\{m\in\mathbb
        N:(a^m)^n=a^{m\cdot n}\}$

        \begin{enumerate}[1.]
          \item $(a^1)^n= a^n=a^{1\cdot n}$. Portanto $1\in S$
          \item Considere $k\in S$, de modo que $(a^k)^n=a^{k\cdot n}$

            Utilizando a propriedade (b):

            $$
            (a^{k+1})^n=(a^k\cdot a)^n=(a^k)^n\cdot a^n
            $$

            E aplicando a hipótese:

            $$
            (a^{s(k)})^n=a^{k\cdot n}\cdot a^n=a^{k\cdot n+n}=a^{n\cdot s(k)}
            $$

            Logo, $s(k)\in S$. E pelo Princípio da Indução Finita a
            propriedade é verdadeira para todos os números naturais.

        \end{enumerate}

      \item $(a\cdot b)^n=a^n \cdot b^n$

        \textbf{Solução:}

        Considere o conjunto $S=\{n\in\mathbb N^*:(a\cdot b)^n=a^n\cdot b^n\}$

        \begin{enumerate}[$1.$]
          \item Para $n=1$, $(a\cdot b)^1=a\cdot b=a^1\cdot b^1$.
            Portanto $1\in S$

          \item Suponha que $k\in S$ tal que $(a\cdot b)^k=a^k\cdot b^k$.

            Multiplicamos ambos os lados por $(a\cdot b)$ e desenvolvendo:

            $$
            (a\cdot b)^k(a\cdot b)=a^k\cdot b^k\cdot a\cdot b
            $$

            $$
            (a\cdot b)^{k+1}=(a^k\cdot a)\cdot
            (b^k\cdot b)=a^{k+1}\cdot b^{k+1}
            $$

            $$
            (a\cdot b)^{s(k)}=(a^k\cdot a)\cdot
            (b^k\cdot b)=a^{s(k)}\cdot b^{s(k)}
            $$

            Portanto $s(k)\in S$. E pelo Princípio da Indução Finita,
            é verdadeiro para todos os números naturais.

        \end{enumerate}
    \end{enumerate}

  \item Mostre as seguintes fórmulas por indução

    \begin{enumerate}[\bfseries (a)]
      \item $1^2+2^2+...+n^2=\frac{n(n+1)(2n+1)}{6}$

        \textbf{Solução:}

        Considere o conjunto $S=\{n\in\mathbb
        N:1^2+2^2+...+n^2=\frac{n(n+1)(2n+1)}{6}\}$

        \begin{enumerate}[$1.$]
          \item Para $n=1$, $\frac{1(1+1)(2\cdot 1+1)}{6}=1$. Portanto $1\in S$

          \item Suponha que $n\in S$ tal que
            $1^2+2^2+...+n^2=\frac{n(n+1)(2n+1)}{6}$.

            Podemos somar $(n+1)^2$ em ambos os lados:

            $$
            1^2+2^2+...+n^2+(n+1)^2=\frac{n(n+1)(2n+1)}{6} + (n+1)^2
            $$

            $$
            =(n+1)(\frac{n(2n+1)}{6} + (n+1))=(n+1)(\frac{2n^2+n+6n+6}{6})
            $$

            $$
            =(n+1)(\frac{2n^2+7n+6}{6})=\frac{(n+1)(2n(n+2)+3(n+2))}{6}
            $$

            $$
            =\frac{(n+1)(n+2)(2(n+1)+1)}{6}=\frac{s(n)(s(n)+1)(2s(n)+1)}{6}
            $$

            Portanto $s(n)\in S$. E pelo Princípio da Indução Finita,
            é verdadeiro para todos os números naturais.

        \end{enumerate}

      \item $\frac{1}{1\cdot
        2}+\frac{1}{2\cdot3}+...+\frac{1}{n(n+1)}=\frac{n}{n+1}$

        \textbf{Solução:}

        Considere o conjunto $S=\{n\in\mathbb
        N:\frac{1}{1\cdot2}+\frac{1}{2\cdot3}+...+\frac{1}{n(n+1)}=\frac{n}{n+1}\}$

        \begin{enumerate}[$1.$]
          \item Para $n=1$, $\frac{1}{1(1+1)}=\frac{1}{1+1}$. Portanto $1\in S$

          \item Suponha que $n\in S$

            $$
            \frac{1}{1\cdot2}+\frac{1}{2\cdot3}+...+\frac{1}{n(n+1)}=\frac{n}{n+1}
            $$

            Adicionando $\frac{1}{(n+1)(n+2)}$ de ambos os lados:

            $$
            =\frac{n}{n+1}+\frac{1}{(n+1)(n+2)}=\frac{n(n+2)}{(n+1)(n+2)}+\frac{1}{(n+1)(n+2)}
            $$

            $$
            =\frac{n^2+2n+1}{(n+1)(n+2)}=\frac{n(n+1)+(n+1)}{(n+1)(n+2)}=\frac{(n+1)(n+1)}{(n+1)(n+2)}=
            $$

            Simplificando:

            $$
            =\frac{n+1}{n+2}=\frac{s(n)}{s(n)+1}
            $$

            Portanto $s(n)\in S$. E pelo Princípio da Indução Finita,
            é verdadeiro para todos os números naturais.

        \end{enumerate}

    \end{enumerate}

  \item Uma progressão aritmética (PA) é uma sequência de números reais
    ($a_n$) tal que $a_1$ é dado e, para todo $n\in\mathbb N$, tem-se

    {\center$a_{n+1}=a_n+r$\endcenter}

    \hbox{onde $r$ é um número real fixo chamado razão}

    \begin{enumerate}[\bfseries (a)]
      \item Mostre que $a_n=a_1+(n-1)r$

        \textbf{Solução:}

        Considere o conjunto $S=\{n\in\mathbb N:a_n=a_1+(n-1)r\}$

        \begin{enumerate}[$1.$]
          \item Para $n=1$, $a_1=a_1+(1-1)r=a_1$. Portanto $1\in S$

          \item Suponha que $n\in S$ de modo que $a_n=a_1+(n-1)r$

            Pela definição:

            $$
            a_{n+1}=a_n+r
            $$

            E usando a hipótese:

            $$
            a_{n+1}=a_1+(n-1)r+r=a_1+nr
            $$

            Portanto $n+1\in S$. E pelo Princípio da Indução Finita,
            é verdadeiro para todos os números naturais.

        \end{enumerate}

      \item Se $a_n=a_1+a_2+...+a_n$. Mostre que
        $S_n=na_1+\frac{n(n-1)r}{2}=\frac{(a_1+a_n)n}{2}$

        \textbf{Solução:}

        Considere o conjunto $S=\{n\in\mathbb
        N:S_n=na_1+\frac{n(n-1)r}{2}\}$

        \begin{enumerate}[$1.$]
          \item Para $n=1$, $S_n=1a_1+\frac{1(1-1)r}{2}=a_1$. Portanto $1\in S$

          \item Suponha que $n\in S$

            Sabendo que $S_{n+1}=S_n+a_{n+1}$, utilizando (a) e
            assumindo a hipótese:

            $$
            S_{n+1}=na_1+\frac{n(n-1)r}{2}+a_{n+1}=na_1+\frac{n(n-1)r}{2}+a_1+nr
            $$

            $$
            =(n+1)a_1+\frac{n^2r-nr+2nr}{2}=(n+1)a_1+\frac{n(n+1)r}{2}=s(n)a_1+\frac{s(n)(s(n)-1)r}{2}
            $$

            Portanto $n+1\in S$. E pelo Princípio da Indução Finita,
            é verdadeiro para todos os números naturais.

            E podemos provar a relação de igualdade também:

            $$
            na_1+\frac{n(n-1)r}{2}=\frac{2na_1+n(n-1)r}{2}
            $$

            $$
            =\frac{n(a_1+a_1+(n-1)r)}{2}=\frac{n(a_1+a_n)}{2}
            $$

        \end{enumerate}
    \end{enumerate}

    \newpage

  \item Prove por indução.

    \begin{enumerate}[\bfseries (a)]
      \item $2^n>n$ para todo $n\in\mathbb N$

        \textbf{Solução:}

        Seja $S=\{n\in\mathbf N:2^n>n\}$.
        \begin{enumerate}[1.]
          \item $2^n = 2 > 1$. Portanto $1\in S$
          \item Suponha que $n\in S$

            Podemos multiplicar ambos os lados por $2$:

            $$
            2^n\cdot2 > 2n
            $$

            $$
            2^{n+1} > 2n
            $$

            Como $n$ é positivo, logo: $2n = n + n > n + 1$

            $$
            2^{n+1} > 2n > n + 1\implies 2^{s(n)}>s(n)
            $$

            E portanto $s(n)\in S$. E pelo Princípio da Indução
            Finita a propriedade é valida para todos os números naturais.

        \end{enumerate}
    \end{enumerate}

\end{enumerate}

\end{document}
